\documentclass[11pt]{article}
\usepackage[T1]{fontenc}
\usepackage[utf8]{inputenc}
\usepackage{lmodern}
\usepackage[english]{babel}
\usepackage{amsmath,amssymb,mathtools,array,booktabs}
\usepackage{geometry}
\usepackage{microtype}
\geometry{a4paper,margin=1.45cm}
\setlength{\parindent}{1.25em}
\setlength{\parskip}{0pt}
\allowdisplaybreaks
\setlength{\emergencystretch}{10em}
\sloppy
\hbadness=10000
\vbadness=10000
\hfuzz=2pt
\vfuzz=10pt
\raggedbottom
\providecommand{\frR}{\mathfrak R}
\providecommand{\frT}{\mathfrak T}
\providecommand{\frM}{\mathfrak M}
\providecommand{\frN}{\mathfrak N}
\providecommand{\frO}{\mathfrak O}
\providecommand{\frS}{\mathfrak S}
\providecommand{\frB}{\mathfrak B}
\providecommand{\frC}{\mathfrak C}
\providecommand{\frA}{\mathfrak A}
\providecommand{\frakp}{\mathfrak p}
\providecommand{\ideal}[1]{\mathfrak{#1}}
\providecommand{\qmod}[1]{\pmod{#1}}

% combined paper 30 macro support
\providecommand{\mR}{\mathfrak{R}}
\providecommand{\mK}{\mathfrak{K}}
\providecommand{\mL}{\mathfrak{L}}
\providecommand{\mS}{\mathfrak{S}}
\providecommand{\mT}{\mathfrak{T}}
\providecommand{\mM}{\mathfrak{M}}
\providecommand{\mN}{\mathfrak{N}}
\providecommand{\mA}{\mathfrak{A}}
\providecommand{\mB}{\mathfrak{B}}
\providecommand{\mC}{\mathfrak{C}}
\providecommand{\mD}{\mathfrak{D}}
\providecommand{\mE}{\mathfrak{E}}
\providecommand{\mF}{\mathfrak{F}}
\providecommand{\ma}{\mathfrak{a}}
\providecommand{\mb}{\mathfrak{b}}
\providecommand{\mc}{\mathfrak{c}}
\providecommand{\md}{\mathfrak{d}}
\providecommand{\me}{\mathfrak{e}}
\providecommand{\mf}{\mathfrak{f}}
\providecommand{\mm}{\mathfrak{m}}
\providecommand{\mn}{\mathfrak{n}}
\providecommand{\mo}{\mathfrak{o}}
\providecommand{\mpideal}{\mathfrak{p}}
\providecommand{\mq}{\mathfrak{q}}
\providecommand{\mt}{\mathfrak{t}}
\providecommand{\mv}{\mathfrak{v}}
\providecommand{\mbar}[1]{\overline{#1}}


% added macros for complete Paper 30 source/control
\providecommand{\mG}{\mathfrak{G}}
\providecommand{\mH}{\mathfrak{H}}
\providecommand{\mU}{\mathfrak{U}}
\providecommand{\mV}{\mathfrak{V}}
\begin{document}


\setcounter{footnote}{25}

\subsection*{§6. Ideal theory under the assumption of the divisor-chain condition}

Let the basis be a commutative ring $\mR$ for which axiom I, the divisor-chain condition, is satisfied. In what follows assume also a fixed well-ordering of all elements of $\mR$. This also determines a well-ordering of all ideals of $\mR$. Indeed, the two assumptions imply at once that every ideal of $\mR$ possesses an ideal basis consisting of finitely many elements. Passing from the well-ordering of the elements of $\mR$ to a quasi-lexicographic well-ordering of all finite subsets of $\mR$,
\footnote{This is to be understood as follows. If the elements $a_1,\ldots,a_n$ of a finite subset $\mathfrak W$ are denoted so that the ordering of the indices agrees with the prescribed well-ordering in $\mR$, then every subset with $n$ elements precedes one with $m>n$ elements. If two subsets have the same number of elements, the first precedes the second when the first element at which their ordered lists differ precedes the corresponding element of the other list. Thus every system of finite subsets has a first element. Given a well-ordering of $\mR$, no further choice postulate is needed; in particular, if $\mR$ is countable, as in the case of an algebraic number field, no choice postulate is involved. The role of the choice postulate in ideal theory was mentioned, without further detail, in \emph{Idealtheorie}, note 5.}
and assigning to each ideal, as distinguished basis, the first among its possible bases in this well-ordering of finite subsets, one obtains a well-ordering of the ideals together with that of the finite subsets.

\textbf{Theorem I.} \emph{Under the assumption of the divisor-chain condition, every ideal of $\mR$ admits a representation as the least common multiple of finitely many irreducible ideals, i.e. of ideals which cannot be represented as the least common multiple of two proper divisors.}

For the proof it is enough to show: if Theorem I is false for an ideal $\mm$, then $\mm$ possesses a proper divisor for which Theorem I is likewise false; from this one can construct, contrary to the divisor-chain condition, a divisor chain not terminating after finitely many steps. Indeed $\mm$ must be reducible, since otherwise $\mm=[\mm]$ would already give the desired representation. If $\mm=[\ma,\mb]$, Theorem I cannot hold for both $\ma$ and $\mb$ at the same time; otherwise the corresponding representation would follow for $\mm$. Hence there are proper divisors of $\mm$ for which Theorem I is false; let $\ma_1$ be the first of these in the well-ordering of ideals. Constructing in the same way, from $\ma_1$, a proper divisor $\ma_2$, and in general from $\ma_i$ a proper divisor $\ma_{i+1}$, one obtains a well-defined divisor chain which does not terminate, contradicting the assumption.
\footnote{Theorem I plainly holds in the same way for module domains, if the divisor-chain condition is assumed for the system of all modules. The theorem has a purely set-theoretic character; it is the only place in the proof where the well-ordering of ideals is used. Abstractly: in a set $M$ let a family $\Sigma$ of subsets be distinguished and well-ordered. Assume the chain condition for $\Sigma$-sets: every chain $X_1,X_2,\ldots$ in which $X_i$ is a proper superset of $X_{i-1}$ terminates. A $\Sigma$-set is reducible if it is the intersection of two $\Sigma$-sets which are both proper supersets, otherwise irreducible. Then every $\Sigma$-set is an intersection of finitely many irreducible $\Sigma$-sets.}

Because the divisor-chain condition is assumed, §5, 2 permits us simply to speak of primary ideals. The connection between primary and irreducible ideals is the following.

\textbf{Theorem II.} \emph{Under the assumption of the divisor-chain condition, every irreducible ideal is primary; equivalently, every non-primary ideal is reducible.}

Passing to a residue-class ring $\mR/\mm$ preserves the divisor-chain condition by homomorphism. The representation of $\mm$ as a least common multiple corresponds to the representation of the zero ideal; and by the first isomorphism theorem (§4, 3), primary divisors of $\mm$ correspond again to primary ideals in $\mR/\mm$ and conversely. Thus it is enough to consider the decomposition of the zero ideal in the residue-class ring. Since that ring is again a ring of the same generality, one may from the outset suppose the ideal to be decomposed to be the zero ideal of $\mR$.

Suppose then that the zero ideal of $\mR$ is not primary. There is at least one pair of ideals $\ma,\mb$ -- and $\mb$ may be assumed principal -- such that
\[
\ma\ne(0),\qquad \mb^x\ne(0)\quad\text{for every }x,
\qquad\text{but}\qquad \ma\mb=(0).
\]
Form the chain of ideal quotients
\[
\ma,\quad \ma:\mb,\quad \ldots,\quad \ma:\mb^\nu,\quad\ldots .
\]
This chain terminates; let, for instance,
\[
\mt=\ma:\mb^{m-1}=\ma:\mb^m=\cdots .
\]
Thus $\mt=\mt:\mb$, or $\mb$ is prime to $\mt$. Further, $\mt$ is a divisor of $\ma$ different from the zero ideal, and $\mb^x$ is different from zero for every exponent $x$. It therefore suffices, for the proof of Theorem II, to establish
\[
        (0)=[\mt,\mb^{m+1}].
\]
By definition,
\[
        \mb^{m-1}\mt\equiv0\pmod{\ma},\qquad\text{hence}\qquad \mb^m\mt=(0),
\]
so it remains only to show
\[
        \mc=[\mt,\mb^{m+1}]\equiv0\pmod{\mb^m\mt}.
\]
This follows from the assumptions that $\mb$ is principal and prime to $\mt$. Every element $c$ of $\mc$, because of its divisibility by $\mb^{m+1}$, has a representation -- where $n$ denotes an integer symbol --
\[
        c=k b^{m+1}+n b^{m+1}=r b^m,
\]
where $r$ is again an element of $\mR$. From $c\equiv0\pmod{\mt}$ follows $r b^m\equiv0\pmod{\mt}$, and hence $r\equiv0\pmod{\mt}$, since $\mt:\mb=\mt$. Therefore $c\equiv0\pmod{\mt\mb^m}$, and Theorem II is proved.

\textbf{Theorem III.} \emph{Under the assumption of the divisor-chain condition, every ideal admits a shortest representation as the least common multiple of finitely many greatest primary components belonging to distinct prime ideals.}

To obtain such a representation, replace, whenever possible, the representation by irreducible ideals furnished by Theorem I -- hence, by Theorem II, by primary ideals -- by a shortest representation. If the primary ideals belonging to the same prime ideal are collected together, §5, 3 gives the asserted representation. The primary components may be called greatest because the least common multiple of any of them is no longer primary.
\footnote{The corresponding prime ideals and the isolated components are uniquely determined. Compare \emph{Idealtheorie}, or W. Krull, ``Ein neuer Beweis für die Hauptsätze der allgemeinen Idealtheorie,'' Math. Ann. 90 (1923), pp. 55--64. In §7 a direct proof of uniqueness is given for the simple special case occurring there. The primary ideals there recognized as unique are isolated components.}

\subsection*{§7. Ideal theory under the assumption of the double-chain condition}

The simplifications, compared with the theory developed so far, rest on the auxiliary results given under 1 and 2.

1. \emph{If, in a commutative ring without zero divisors, the multiple-chain condition is satisfied, then the ring is already a field.} It is to be shown that, for $a\ne0$, the equation $ax=b$ always has a solution in the ring. That it can have at most one solution follows from the absence of zero divisors.

Let $\ma$ be the principal ideal generated by $a$. By assumption the sequence of powers terminates; let $\ma^m$ be equal to all following powers. If $\mb$ denotes the principal ideal generated by $b$, then
\[
        \ma^m\mb=\ma^{m+1}\mb,
        \qquad \ma^{m+1}\mb=(a^{m+1}b).
\]
In particular the element $a^m b$ has a representation -- again with $n$ an integer symbol --
\[
        a^m b=k a^{m+1}b+n a^{m+1}b=a^{m+1}c,
\]
where $c$ is an element of $\mR$. Since no zero divisors exist, this gives $b=ac$, proving that the ring is a field.

2. From 1 follows at once the auxiliary result: \emph{If the multiple-chain condition is satisfied in a commutative ring, then a prime ideal has no proper divisor different from $\mo$.}

Likewise one obtains the supplement: \emph{If only axiom II is satisfied, that is, the multiple-chain condition modulo every ideal different from the zero ideal, then every prime ideal different from the zero ideal has no proper divisor different from $\mo$.}

For the residue-class ring modulo a prime ideal $\mpideal$ is, by definition, a ring without zero divisors; since the multiple-chain condition is also preserved there by homomorphism, it is a field by 1. From the one-to-one correspondence between the divisors of $\mpideal$ and the ideals in the residue-class ring, $\mpideal$ has no proper divisor other than $\mo$.
\footnote{No identity element need exist in $\mR$, nor therefore in $\mo$, the ideal consisting of all elements of $\mR$, although an identity class exists in the residue-class ring by 1. The simplest example of this kind, in the weaker case of the supplement, is furnished by the system of all even integers.}

\textbf{Theorem IV.} \emph{If the double-chain condition is satisfied in a commutative ring, then in every shortest representation of an ideal the primary components not belonging to the unit ideal $\mo$ are uniquely determined.}

\emph{Supplement.} Under the assumptions of axioms I and II, Theorem IV holds for every ideal different from the zero ideal.

Let
\[
        \mm=[\mq,\mq_1,\ldots,\mq_r]=[\mbar\mq,\mbar\mq_1,\ldots,\mbar\mq_s]
\]
be shortest representations of $\mm$ by greatest primary components, with associated prime ideals
\[
        \mo,\mpideal_1,\ldots,\mpideal_r,
        \qquad \mo,\mbar\mpideal_1,\ldots,\mbar\mpideal_s.
\]
The component $\mq$, respectively $\mbar\mq$, is to be omitted when no primary ideal belonging to $\mo$ occurs. Since
\[
        \mo\,\mpideal_1\cdots\mpideal_r\equiv0\pmod{\mm},
\]
each $\mpideal_i$ must be contained in some $\mbar\mpideal_j$, hence by the auxiliary result is identical with it. Conversely each $\mbar\mpideal_j$ must coincide with some $\mpideal_i$; the prime ideals different from $\mo$ therefore agree. Further,
\[
        \mq\,\mq_1\cdots\mq_r\equiv0\pmod{\mq_i}
\]
implies $\mbar\mq_i\equiv0\pmod{\mq_i}$, because the remaining components are not divisible by $\mpideal_i$ and hence are prime to $\mq_i$. The converse congruence follows in the same way, so the non-$\mo$ primary components are unique. If $\mq$ actually occurs, then $\mbar\mq$ must actually occur as well, because the representations are shortest; this also gives the uniqueness of the associated prime ideals. The proof also shows that every representation without a primary component belonging to $\mo$ is already shortest.

3. If to the double-chain condition one adds the existence of an identity element, Theorem IV sharpens to the following.

\textbf{Theorem V.} \emph{In a commutative ring with identity in which the double-chain condition is satisfied, every ideal can be represented uniquely as a product of finitely many pairwise coprime primary ideals.}

\emph{Supplement.} Under axioms I, II, III, Theorem V holds for every ideal different from the zero ideal.

Because an identity element exists, $\mo^2=\mo$; hence every primary ideal belonging to $\mo$ is equal to $\mo$ and cannot occur in a shortest representation of an ideal different from $\mo$. Thus, by Theorem IV, the primary components are uniquely determined, and at the same time every representation becomes shortest after omitting $\mo$. By the auxiliary result under 2, two distinct prime ideals are always coprime; by the calculation rules of §4, 4, the same holds for their primary components. The least common multiple therefore becomes their product.

From the same assumptions follows the auxiliary result: \emph{In a commutative ring with identity and the double-chain condition, the prime ideals prime to an ideal $\mm$ are exactly those, apart from the unit ideal, which are not contained in $\mm$. The notions prime and coprime coincide.}

\emph{Supplement.} Under axioms I, II, III, the ideal $\mm$ must be assumed different from the zero ideal.

If $\mpideal$ is not contained in $\mm$, then it is different from all prime ideals belonging to $\mm$, and by the auxiliary result under 2 it is not divisible by any of them. Hence it is prime to each primary component and therefore to $\mm$; at the same time it is coprime to each component and hence to $\mm$.

If, however, $\mpideal\ne\mo$ is contained in $\mm$, then it must coincide with one of the associated prime ideals. Let it belong to the component $\mq=\mq_1$. Then $\mq:\mpideal$ is a proper divisor of $\mq$, since
\[
        \mpideal^{\varrho-1}\equiv0\pmod{\mq:\mpideal},\qquad
        \mpideal^{\varrho-1}\not\equiv0\pmod{\mq}.
\]
At the same time $\mq:\mpideal$, as a divisor of $\mpideal^{\varrho-1}$, is divisible only by the prime ideal $\mpideal\ne\mo$ and is therefore primary. By the uniqueness of the decomposition, the product
\[
        (\mq:\mpideal)\mq_2\cdots\mq_r,
\]
which is divisible by $\mm:\mpideal$, is a proper divisor of $\mm$; hence $\mpideal$ is not prime to $\mm$.

Thus an ideal $\mt$ is prime to $\mm$ if and only if none of the prime ideals belonging to $\mt$ is contained in $\mm$; but in that case $\mt$ is also coprime to $\mm$. Conversely, coprime ideals are always mutually prime (§5, 5), so the two concepts coincide under these assumptions.

\subsection*{§8. Ideal theory under integral closedness in the quotient field}

The ideal theory developed so far is now to be sharpened to the usual one by adjoining the last two axioms.

1. \emph{Auxiliary result.} Let $\mR$ be a commutative ring without zero divisors and with an identity element. If the divisor-chain condition is assumed in $\mR$ (axioms I, III, IV), then from $\ma=\ma\mb$ and $\ma\ne(0)$ it follows that $\mb=\mo$. Hence
\[
        \mc\ne\mc\mt
\]
for all ideals different from the zero and unit ideals.
\footnote{On the other hand, under the same assumptions one cannot infer $\mb=\mc$ from $\ma\mb=\ma\mc$. For example, let $\mR$ be the polynomial domain in $x,y$ over a field, with $\ma=(xy)$, $\mb=(x^3,xy,y^2)$, $\mc=(x^2,y^2)$. Then $\ma\mb=\ma\mc=\ma^2$, but $\mb\ne\mc$.}
The proof is the same as for the auxiliary result §1, 3, since $\ma\mb$ may be regarded as a finite module over $\mb$. Let $\alpha_1,\ldots,\alpha_n$ be an ideal basis of $\ma$ existing by I. From $\ma=\ma\mb$ one obtains the system
\[
        \alpha_i=b_{i1}\alpha_1+\cdots+b_{in}\alpha_n,
        \qquad b_{ij}\equiv0\pmod{\mb},\quad i=1,\ldots,n .
\]
Since $\mR$ has no zero divisors, the determinant $|b_{ij}-e_{ij}|$ vanishes; with $e_{ij}=0$ for $i\ne j$ and $e_{ii}=e$. From
\[
        e-b_{11}-\cdots\equiv0\pmod{\mb}
\]
it follows that $e\equiv0\pmod{\mb}$ and therefore $\mb=\mo$.

2. It remains to show only that, after adjoining integral closedness in the quotient field (axiom V) to axioms I through IV, every primary ideal different from the zero ideal is a power of its associated prime ideal. The uniqueness of the resulting representation
\[
        \mm=\mpideal_1^{\varrho_1}\cdots\mpideal_r^{\varrho_r}
\]
for every ideal different from the zero ideal has already been proved by Theorem V and the auxiliary result under 1. At the same time it has been shown that these prime ideals have no proper divisor distinct from the unit ideal. The proof is first given for exponent two, using Dedekind's corollary II (§1, 4), and then in general by complete induction.

\emph{Auxiliary result.} \emph{Under axioms I through V there are no primary ideals of exponent two other than $\mq=\mpideal^2$.}

It is to be shown that from
\[
        \mpideal^2\equiv0\pmod{\mq},
        \qquad \mq\equiv0\pmod{\mpideal},
        \qquad \mq\not\equiv0\pmod{\mpideal^2}
\]
there necessarily follows $\mq=\mpideal^2$. Choose
\[
        \mc\equiv0\pmod{\mq},
        \qquad \mc\not\equiv0\pmod{\mpideal},
        \qquad \mc\equiv0\pmod{\mpideal^2}.
\]
From the auxiliary result §7, 3, it follows that $\mc:\mpideal$ is a proper divisor of $\mc$; hence there is an element $b$ such that
\[
        b\mpideal\equiv0\pmod{\mc}\equiv0\pmod{\mq},
        \qquad b\not\equiv0\pmod{\mc}.
\]
Thus $y=b/c$ belongs to the quotient field but not to $\mR$; it is not integral. We must prove $b\not\equiv0\pmod{\mpideal}$; then, since $\mq$ is primary and belongs to $\mpideal$, it follows that $\mpideal\equiv0\pmod{\mq}$.

By Dedekind's corollary II there exist elements $m,n$ of $\mR$ such that
\[
        y=b/c=m/n,
        \qquad m^2/n \text{ is not integral};
\]
that is,
\[
        bn=mc,
        \qquad m\not\equiv0\pmod{\mn},
        \qquad m^2\not\equiv0\pmod{\mn}.
\]
Multiplying $b\mpideal$ by $m$ gives
\[
        mb\mpideal\equiv0\pmod{\mn c},
        \qquad\text{hence}\qquad m\mpideal\equiv0\pmod{\mn}.
\]
Therefore $m\not\equiv0\pmod{\mpideal}$, because $m^2\not\equiv0\pmod{\mn}$ and $\mn\not\equiv0\pmod{\mpideal}$; the latter again follows from §7, 3, since $\mn:\mpideal$ is a proper divisor of $\mn$. The four relations
\[
        m\not\equiv0\pmod{\mpideal},\qquad
        n\not\equiv0\pmod{\mpideal},\qquad
        c\equiv0\pmod{\mpideal},\qquad
        c\not\equiv0\pmod{\mpideal^2},\qquad
        bn=mc
\]
now imply $b\not\equiv0\pmod{\mpideal}$. For, since $\mpideal^2$ is primary, one first gets $mc\not\equiv0\pmod{\mpideal^2}$, and then $b\not\equiv0\pmod{\mpideal}$ from $bn=mc$. This proves the auxiliary result.

3. \emph{Auxiliary result.} \emph{Under axioms I through V there are no primary ideals of exponent $\varrho$ other than $\mq=\mpideal^\varrho$.}

This follows from the auxiliary result under 2 without further use of the axioms.
\footnote{The proof of auxiliary result 3 under the assumption of auxiliary result 2 is found in Masazo Sono, \emph{On Congruences II}, §§9 and 10. Compare note 17.}
First one proves:
\[
        \text{if }\mc\equiv0\pmod{\mpideal},\ \mc\not\equiv0\pmod{\mpideal^2},
        \quad\text{then}\quad
        \mpideal^\varrho=(\mc^\varrho,\mpideal^{\varrho+1})
\]
for every $\varrho$. By the result under 2, $\mpideal=(\mc,\mpideal^2)$. Assuming $\mpideal^{\varrho-1}=(\mc^{\varrho-1},\mpideal^\varrho)$, multiplication by $\mpideal$ and the distributive law give
\[
        \mpideal^\varrho=(\mc^\varrho,\mpideal^{\varrho+1}).
\]
The desired result may be put in the following second form: if
\[
        \mq\equiv0\pmod{\mpideal^\varrho},
        \qquad \mq\not\equiv0\pmod{\mpideal^{\varrho+1}},
        \qquad \mpideal^{\varrho+1}\equiv0\pmod{\mq},
        \quad \varrho\ge1,
\]
then already $\mpideal^\varrho\equiv0\pmod{\mq}$. Indeed, for every primary ideal there is such an exponent $\varrho>1$; the condition $\mq\not\equiv0\pmod{\mpideal^{\varrho+1}}$ follows from $\mpideal^\varrho\equiv0\pmod{\mq}$ and $\mpideal^\varrho\ne\mpideal^{\varrho+1}$ by the result under 1. Repeated application of the second form gives $\mpideal^\varrho\equiv0\pmod{\mq}$ and hence $\mq=\mpideal^\varrho$.

From the assumption of the second form and from the representation of $\mpideal^\varrho$, each element $q$ of $\mq$ can be written
\[
        q=b c^\varrho+p^{\varrho+1},
        \qquad b\equiv0\pmod{\mpideal},
        \quad\text{or}\quad b c^\varrho\equiv0\pmod{\mq,\mpideal^{\varrho+1}}.
\]
Since $(\mq,\mpideal^{\varrho+1})$ is primary, it follows that
\[
        c^\varrho\equiv0\pmod{\mq,\mpideal^{\varrho+1}},
        \qquad\text{or}\qquad
        c^{\varrho+1}\equiv0\pmod{\mq,\mpideal^{\varrho+2}},
\]
and hence $\mpideal^\varrho\equiv0\pmod{\mq}$.

4. Summarizing, one obtains:

\textbf{Theorem VI.} \emph{If axioms I through V hold in a commutative ring, then every ideal different from the zero and unit ideals can be represented uniquely as a product of powers of finitely many prime ideals different from the zero and unit ideals. These prime ideals have no proper divisor different from the unit ideal.}

\subsection*{§9. The axioms as consequences of the assumed decomposition}

Conversely, in order to infer the axioms from the existence of the usual ideal decomposition -- which by Theorem VI follows from axioms I through V -- the decomposition must be assumed in the following form.

\emph{Assumption.} In the commutative ring $\mR$, every ideal not generated by the zero or the unit element is representable uniquely as a product of powers of prime ideals. These prime ideals are simple ideals not generated by the unit element, i.e. they have no proper divisor different from $\mo$; conversely all simple ideals are prime ideals. The uniqueness is to hold in the sharp form that
\[
        \ma^\varrho\ne\ma^{\varrho+1}
\]
for every ideal not generated by the zero or by the unit element.

The existence of an identity element is not assumed here. If no identity element exists, the condition ``not generated by the unit element'' is no restriction.

1. \emph{Proof of axiom III, the existence of the identity element.} Suppose axiom III is not satisfied. It is to be shown that $\mo^2$ becomes a simple ideal without being prime, contrary to the assumption. By the sharp uniqueness assumption, $\mo\ne\mo^2$; hence $\mo\equiv0\pmod{\mo^2}$, but $\mo^2\not\equiv0\pmod{\mo^2}$, and so $\mo^2$ is not prime. It is, however, simple: it cannot be divisible by any prime ideal different from $\mo$, and by the assumed product representation it has no divisors except $\mo$ and $\mo^2$.

2. \emph{Proof of axiom IV, absence of zero divisors.} If $a$ is a zero divisor, say $ab=0$ with $a\ne0$ and $b\ne0$, then multiplication of the product representations of the principal ideals generated by $a$ and $b$ gives
\[
        (0)=\mpideal_1^{\varrho_1}\cdots\mpideal_r^{\varrho_r},
\]
where the $\mpideal_i$ are different from the zero and unit ideals, the latter because the identity element has already been proved to exist. Take the representation to be shortest, in the sense that deleting any factor yields a proper divisor of the zero ideal; this need not be automatic, since no uniqueness assumption has been made about the zero ideal. If only one prime ideal occurs, then $(0)=\mpideal^\varrho=\mpideal^{\varrho+1}$, contrary to the sharp uniqueness demand. If more than one prime ideal occurs, then
\[
\begin{aligned}
\mpideal_1^{\varrho_1}
 &= (\mpideal_1^{\varrho_1},\mpideal_2^{\varrho_2}\cdots\mpideal_r^{\varrho_r})\mpideal_1^{\varrho_1}  \\
 &= \mpideal_1^{\varrho_1}\mpideal_2^{\varrho_2}\cdots\mpideal_r^{\varrho_r},
\end{aligned}
\]
because, by the existence of the identity element, $\mpideal_1$ is coprime to all the other prime ideals. This again contradicts the sharp uniqueness condition.

3. \emph{Proof of axioms I and II, the chain conditions.} The assumptions give, for every ideal different from the zero ideal: divisibility is reflected in the product representation. Thus from $\mm\equiv0\pmod{\ma}$ and $\ma\ne\mo$, with
\[
        \mm=\mpideal_1^{\varrho_1}\cdots\mpideal_r^{\varrho_r},
        \qquad
        \ma=\mbar\mpideal_1^{\bar\varrho_1}\cdots\mbar\mpideal_s^{\bar\varrho_s},
\]
it follows that every $\mbar\mpideal_j$ is identical with some $\mpideal_i$ and that $0\le\bar\varrho_i\le\varrho_i$. Hence there are only finitely many divisors of $\mm$, corresponding to the finitely many exponent combinations. Therefore the double-chain condition holds in the residue-class ring modulo $\mm$. Axioms I and II are thus proved: the divisor-chain condition also holds in $\mR$ itself, since every proper divisor of the zero ideal is different from the zero ideal.

The proof of axiom V, integral closedness in the quotient field, uses the standard consequences of ideal decomposition, which must therefore first be derived.

4. \emph{The residue-class ring modulo every ideal different from the zero ideal is a principal-ideal ring.} The ideal may be assumed different from the unit ideal, since for the residue-class ring consisting only of the zero element the assertion is immediate.

If the ideal is first primary, $\mq=\mpideal^\varrho$, and if $\mc\equiv0\pmod{\mpideal}$ but $\mc\not\equiv0\pmod{\mpideal^2}$, then by 3 one has $\mc=\mpideal\ma$ with $(\ma,\mpideal)=\mo$. Consequently
\[
        \mpideal^\sigma=(\mc^\sigma,\mpideal^{\sigma+1})
        \qquad(\sigma<\varrho),
\]
so every ideal in the residue-class ring modulo a primary ideal is principal. In general, if
\[
        \mm=\mq_1\cdots\mq_r
\]
is the representation and
\[
        \mR/\mm=\ma_1+\cdots+\ma_r
\]
the corresponding representation of the residue-class ring as a direct sum (§4, 5), then $\ma_i$ is isomorphic to $\mR/\mq_i$; every ideal of $\ma_i$ is therefore principal. If $\mc$ is any ideal of the residue-class ring, then $\ma_i\mc$ is a principal ideal $\ma_i c_i$, and
\[
        \mc=\ma_1c_1+\cdots+\ma_r c_r=\mo c,
        \qquad c=c_1+\cdots+c_r.
\]
Indeed $\ma_i\mo c_i$ and $\ma_i\mc$ agree in $\ma_i$, since $\ma_i\ma_k=(0)$ for $i\ne k$ and $c_i\equiv0\pmod{\ma_i}$.

The theorem on principal-ideal residue rings also has the following form. If $\mc$ is any ideal, it can be transformed into a principal ideal by multiplication with an ideal prime to a prescribed ideal $\mb$. For, putting $\mm=\mb\mc$, the ideal $\mc$ becomes principal modulo $\mm$; that is,
\[
        \mc=(\mo c,\mm)=(\mc\ma,\mc\mb)=\mc(\ma,\mb),
\]
so that $\mo c=\mc\ma$ and $(\ma,\mb)=\mo$.

5. \emph{Theory of fractional ideals.} A fractional ideal means a finite $\mR$-module in the quotient field $\mK$.

\emph{The nonzero finite $\mR$-modules in $\mK$ form an abelian group under multiplication.}

The product of two finite $\mR$-modules is again finite; multiplication is associative and commutative; since $\mK=\mo\mK$, the unit ideal is the identity element. It remains only to show that the equation
\[
        \mathfrak U X=\mo
\]
always has a solution in the system of finite $\mR$-modules.

Preliminary remark: if the equation $\mathfrak A T=\mathfrak B$ has one and only one solution, then $T$ is the module quotient $\mathfrak B:\mathfrak A$ (§5, 4). For
\[
        T\equiv0\pmod{\mathfrak B:\mathfrak A},
        \qquad
        \mathfrak B=\mathfrak A T\equiv0\pmod{\mathfrak A(\mathfrak B:\mathfrak A)}\equiv0\pmod{\mathfrak B}.
\]

If first $\mathfrak A$ is principal, $\mathfrak A=(\alpha)$, then $X=(e/\alpha)$ is a solution of $\mathfrak A X=\mo$, and $X$ is again a finite $\mR$-module. It follows that every finite $\mR$-module is the module quotient of two ideals of $\mR$, which justifies the term fractional ideal. Indeed, if $t_1,\ldots,t_n$ is a module basis of $T$, put $t_i=\tau_i/a$ and let $\mt$ be the ideal generated by the $\tau_i$ in $\mR$. Then $\mo a\,T=\mt$, whence by the preliminary remark $T=(\mt:\mo a)$, the quotient being taken in $\mK$.

The solution of $\mathfrak A X=\mo$ can now be given generally. Put $\mathfrak A=\ma:\mo c$, and choose $\mb$ so that $\ma\mb$ is the principal ideal $\mo a$. Then
\[
        \mathfrak A\,\mo c=\ma,
        \qquad \mathfrak A\,\mb\mo c=\mo a,
        \qquad \mathfrak A(\mb\mo c:\mo a)=\mo.
\]
Here $X$, as the quotient of an ideal by a principal ideal, is again a finite $\mR$-module. The group property is proved. It also follows that the module quotient of any two ideals, $\ma:\mb$, is a finite $\mR$-module, as the solution of $\mb X=\ma$.

6. From the group property follow further consequences.

6\,$\alpha$. \emph{Cancellation is possible: from $\mathfrak T=\mathfrak A\mathfrak M:\mathfrak B\mathfrak M$ follows $\mathfrak T=\mathfrak A:\mathfrak B$, and conversely.} For $\mathfrak T\mathfrak B\mathfrak M=\mathfrak A\mathfrak M$ gives $\mathfrak T\mathfrak B=\mathfrak A$, and conversely.

6\,$\beta$. \emph{Every fractional ideal admits a representation $\mathfrak C=\ma:\mb$, where $\ma$ and $\mb$ are coprime; this condition determines $\ma$ and $\mb$ uniquely.} Existence follows from 5 and 6\,$\alpha$. From $\mathfrak C=\ma:\mb=\mbar\ma:\mbar\mb$ follows $\ma\mbar\mb=\mb\mbar\ma$, and hence uniqueness when both pairs $\ma,\mb$ and $\mbar\ma,\mbar\mb$ are assumed coprime.

6\,$\gamma$. \emph{Every principal module generated by a nonintegral element, i.e. by an element of $\mK$ not in $\mR$, admits a representation as a quotient of principal ideals such that no power of the numerator is divisible by the denominator.} Let the reduced representation be $\mo n=\mb:\mc$, so $(\mb,\mc)=\mo$; choose $\ma$ according to 4 so that $\mo b=\mb\ma$ and $(\ma,\mc)=\mo$. Then
\[
        \mo n=\ma\mb:\ma\mc=\mo b:\ma\mc,
\]
and since $\ma\mc=\mo b:\mo n$, the ideal $\ma\mc$ is principal as well; write
\[
        \mo n=\mo b:\mo c.
\]
No power of $b$ is divisible by $c$, because $(\ma,\mc)=\mo$ and $(\mb,\mc)=\mo$.

7. \emph{Proof of axiom V, integral closedness in the quotient field.} By 6\,$\gamma$, every nonintegral element of $\mK$ has a quotient representation
\[
        \eta=a/c
\]
such that, in $\mR$, no power of the numerator is divisible by the denominator. Such an element cannot satisfy an equation characterizing it as integral over $\mR$; for from
\[
        \eta^n+r_1\eta^{n-1}+\cdots+r_n=0
\]
one would obtain $a^n\equiv0\pmod{\mo c}$ in $\mR$.

8. \emph{Consequence.} If axioms I through IV hold in a commutative ring $\mR$, and every primary ideal is irreducible, then axiom V, integral closedness in the quotient field, also holds. Because of axioms I through III, every prime-ideal power $\mpideal^\varrho$ is primary; in particular $\mpideal^2$ is primary and therefore irreducible by assumption. From this one must prove
\[
        \mpideal=(\mo c,\mpideal^2).
\]
Then, by the auxiliary result §8, 3, every primary ideal is a prime-ideal power; together with the other assumptions, 7 gives integral closedness.

By the divisor-chain condition,
\[
        \mpideal=(\mo c_1,\ldots,\mo c_n),
\]
where only residue classes modulo $\mpideal$ are relevant as multipliers of the $c_i$, and the $c_i$ may be assumed linearly independent over the residue-class field modulo $\mpideal$. If $n>1$, this linear independence gives a representation
\[
        \mpideal^2=[\mq_1,\mq_2],
        \qquad
        \mq_1=(\mo c_1,\mpideal^2),
        \qquad
        \mq_2=(\mo c_2,\ldots,\mo c_n,\mpideal^2),
\]
with both $\mq_1$ and $\mq_2$ proper divisors of $\mpideal^2$; hence $\mpideal^2$ would be reducible, contrary to the assumption. Thus $\mpideal=(\mo c,\mpideal^2)$, and the asserted consequence is proved.

Conversely, as a consequence of axioms I through V, every primary ideal is immediately irreducible: a prime-ideal power $\mpideal^\varrho$ cannot be the least common multiple of proper divisors of the form $\mpideal^\sigma$ and $\mpideal^\tau$.

9. If zero divisors are allowed in the ring, then axioms I through III and integral closedness in the quotient ring do not imply that every primary ideal is irreducible.
\footnote{Compare note 18.}
Let $\mT$ be a ring in which all axioms I through IV except integral closedness hold; such rings exist, for instance the finite orders distinct from the full system of integral quantities in a finite extension field of the quotient field of $\mR$, when axioms I through V hold in $\mR$ (§3, 2). By the consequence under 8 there is in $\mT$ a reducible primary ideal $\mq$. Let $\mpideal^\varrho$ be a proper multiple of $\mq$, and let $\mR$ be the residue-class ring $\mT/\mpideal^\varrho$, in which the ideal corresponding to $\mq$ is a nonzero reducible primary ideal. In $\mR$ axioms I through III hold, and integral closedness in the quotient ring also holds. For every element of $\mT$ not divisible by $\mpideal$ becomes coprime to $\mpideal^\varrho$; hence all regular elements of $\mR$ are units, and $\mR$ coincides with its quotient ring.

\subsection*{§10. The double-chain condition and composition series}

We now show that, for arbitrary module domains (§2) assumed to be well-ordered, the validity of the double-chain condition -- every divisor chain and every multiple chain of modules terminates after finitely many steps -- is equivalent to the existence of a composition series.
\footnote{The modules of the domain are abelian groups with respect to addition, in fact generalized abelian groups, since the multiplier domain consists of the elements of $\mR$. Since they are abelian groups, the concept of composition series coincides here with that of principal series. The theorems of this paragraph remain valid if, under ``modules,'' one understands the totality of normal divisors of an arbitrary noncommutative, even generalized, group. The notation, slightly different from the earlier one, is meant to recall group theory.}
Specializing the module domain to a commutative ring gives the corresponding facts for the system of all ideals of the ring; in 2, throughout, module isomorphism is then to be replaced by ring isomorphism.

A module domain $\mG$ is called \emph{simple} if in $\mG$ there are no $\mR$-modules besides $\mG$ itself and the zero module $\mE$. A module $\mA$ is called \emph{simple in $\mG$} if $\mG/\mA$ is simple.

A multiple chain
\[
        \mG,\mU_1,\ldots,\mU_r,\mE
\]
is called a composition series of $\mG$ of length $r$ if all modules in the chain are distinct and if each module is simple in the preceding one; equivalently, if the residue-class modules, the quotient groups, $\mU_{i-1}/\mU_i$ are simple module domains, as are $\mG/\mU_1$ and $\mU_r/\mE$. By forming the residue-class module $\mG/\mF$, the case in which a composition series exists from $\mG$ down to $\mF\ne\mG$ is reduced to the absolute case.

1. \emph{If the double-chain condition is assumed in $\mG$, then a composition series exists in $\mG$}, provided $\mG\ne\mE$. From the assumed well-ordering of $\mG$ and the divisor-chain condition, one obtains as in §6, by quasi-lexicographic ordering, a well-ordering of the system of all modules. This well-ordering, together with the divisor-chain condition, gives the existence of at least one module simple in $\mG$. Namely, let $\mG_1$ be the first proper divisor of $\mG$ in the well-ordering; generally let $\mG_i$ be the first proper divisor of $\mG_{i-1}$. The well-determined chain
\[
        \mG,\mG_1,\ldots,\mG_i,\ldots
\]
terminates after finitely many steps and therefore leads necessarily to a module $\mA$ simple in $\mG$. If $\mU_1\ne\mE$, the same procedure gives a module $\mU_2$ simple in $\mU_1$; in general, if $\mU_{i-1}\ne\mE$, there is a module $\mU_i$ simple in $\mU_{i-1}$. This gives a well-determined multiple chain
\[
        \mG,\mU_1,\ldots,\mU_i,
\]
which terminates by the multiple-chain condition and is therefore a composition series of $\mG$.

2. \emph{If a composition series exists in $\mG$, then the double-chain condition holds in $\mG$.} The proof is by complete induction and does not require a well-ordering of $\mG$. In the course of the induction one obtains at the same time a proof of the Jordan--Hölder theorem under the sole assumption that a composition series exists.
\footnote{Compare Masazo Sono (note 21), \emph{On Congruences I}, §§11--13, for the case of rings. There the analogue of a composition series in noncommutative groups is also treated: series $\mR,\mU_1,\ldots,\mU_r,\mE$, where $\mU_i$ is an ideal and simple in $\mU_{i-1}$, without necessarily being an ideal in $\mR$. The present arguments remain valid for noncommutative groups if, by a multiple chain, one means a chain in which each $\mU_i$ is normal in $\mU_{i-1}$, and divisor chains are defined correspondingly. Compare also Dedekind, ``Über die von drei Moduln erzeugte Dualgruppe,'' Math. Ann. 53 (1900), pp. 371--403. There, for much more general domains, the Jordan--Hölder theorem is likewise proved under the existence assumption alone; in place of the second isomorphism theorem there occurs a somewhat weaker correspondence relation, and consequently the statement of the theorem is also somewhat weaker. The principle of the proof is identical with the one given here.}

If $\mG$ is simple, so that $r=0$ and the only composition series is $\mG,\mE$, then the following assertions hold.

$\alpha$) Through every module simple in $\mG$ one can draw a composition series.

$\beta$) Jordan--Hölder theorem: every composition series has the same length and the system of quotient groups (residue-class modules) agrees up to order; corresponding quotient groups are isomorphic.

$\gamma$) Through every module different from $\mG$ one can draw a composition series.

$\delta$) The multiple-chain condition holds.

$\varepsilon$) The divisor-chain condition holds.

Assume therefore that $\alpha$)--$\varepsilon$) have been proved for every module domain possessing a composition series of length $r<r+1$. We prove them under the assumption that in $\mG$ there is a composition series
\[
        \mG,\mA,\mA_1,\ldots,\mA_r,\mE
\]
of length $r+1$.

2$\alpha$. If there is no module simple in $\mG$ other than $\mA$, there is nothing to prove. Let $\mB$ be simple in $\mG$ and $\mB\ne\mA$. Since $\mA$ and $\mB$ are both simple in $\mG$ and are distinct, necessarily $(\mA,\mB)=\mG$. Put $\mM=[\mA,\mB]$. By the second isomorphism theorem (§4, 2),
\[
        \mG/\mB\cong \mA/\mM,
        \qquad
        \mG/\mA\cong \mB/\mM.
\]
Thus $\mM$ is simple in $\mA$ and in $\mB$. Since $\mA$ has a composition series of length $r<r+1$, the induction hypotheses $\alpha$ and $\delta$ give in $\mA$ a composition series through $\mM$, say
\[
        \mA,\mM,\mM_1,\ldots,\mM_{r-1},\mE.
\]
Consequently
\[
        \mG,\mB,\mM,\mM_1,\ldots,\mM_{r-1},\mE
\]
is a composition series of $\mG$ passing through $\mB$.

2$\beta$. To prove the Jordan--Hölder theorem, let
\[
\begin{aligned}
&(1)\quad \mG,\mA,\mA_1,\ldots,\mA_r,\mE,\qquad\text{and}\qquad
(2)\quad \mG,\mB,\mB_1,\ldots,\mB_s,\mE
\end{aligned}
\]
be two composition series of $\mG$. If $\mA=\mB$, the result follows from the induction assumption for length $r<r+1$. Otherwise compare with the series constructed under 2$\alpha$,
\[
        (3)\quad \mG,\mB,\mM,\mM_1,\ldots,\mM_{r-1},\mE,
        \qquad
        (4)\quad \mG,\mA,\mM,\mM_1,\ldots,\mM_{r-1},\mE.
\]
By the induction hypothesis, the system of quotient groups of (1) and (3) agrees; likewise that of (2) and (4), since (4) is a composition series of length $r$ through $\mB$ after passing to the appropriate quotient. Thus $s=r$. The isomorphism under 2$\alpha$ gives agreement between (3) and (4), and hence between (1) and (2). This proves the Jordan--Hölder theorem.

2$\gamma$. Preliminary remark. Let $\mA,\mB,\mH$ be modules of $\mG$, and suppose $\mB$ is simple in $\mA$. Then the modules $(\mB,\mH)$ and $(\mA,\mH)$ either coincide, or $(\mB,\mH)$ is simple in $(\mA,\mH)$. By the second isomorphism theorem,
\[
        (\mA,\mB,\mH)/(\mB,\mH)
        \cong
        \mA/[\mA,(\mB,\mH)].
\]
Since $\mB\equiv0\pmod{\mA}$, this is
\[
        (\mA,\mH)/(\mB,\mH)\cong \mA/\mC,
        \qquad \mC=[\mA,(\mB,\mH)].
\]
Here $\mC\equiv0\pmod{\mA}$; on the other hand $\mB\equiv0\pmod{\mC}$, because $\mB\equiv0\pmod{\mA}$ and $\mB\equiv0\pmod{(\mB,\mH)}$. Hence $\mC$ lies between $\mA$ and $\mB$ and is therefore, by assumption, identical with $\mA$ or $\mB$. This proves the preliminary remark.

Now let
\[
        \mG,\mU,\mU_1,\ldots,\mU_r,\mE
\]
be a composition series of $\mG$, and let $\mH$ be an arbitrary module of $\mG$ different from $\mG$. To show that a composition series passes through $\mH$, form
\[
        \mG,(\mU,\mH),(\mU_1,\mH),\ldots,(\mU_r,\mH),\mH.
\]
By the preliminary remark, the distinct modules among these form a composition series from $\mG$ to $\mH$. Thus the first proper term is simple in $\mG$ (this is the special case 2$\alpha$ when it coincides with $\mB$). By 2$\alpha$ there is in $\mG$ a composition series of length $r<r+1$ through that term, and consequently, by the induction hypothesis, a composition series through $\mH$; adjoining $\mG$ gives a composition series of $\mG$ through $\mH$. Repeating the argument shows that such a series may in particular be chosen as an extension of the series just constructed from $\mG$ to $\mH$.

2$\delta$. Proof of the multiple-chain condition. Suppose again that a composition series of length $r+1$ exists in $\mG$, and let
\[
        \mG,\mB,\mB_1,\ldots,\mB_i,\ldots
\]
be a multiple chain, each term being a proper multiple of the preceding one. That the chain begins with $\mG$ is no restriction of generality, since $\mG$ may always be adjoined as first term. By 2$\gamma$ there is a composition series through $\mB$, and hence $\mB$ has a composition series of shorter length. By the induction hypothesis, the chain
\[
        \mB,\mB_1,\ldots,\mB_i,
\]
terminates after finitely many steps. Hence the same is true after adjoining $\mG$.

2$\varepsilon$. Proof of the divisor-chain condition. Again assume a composition series of length $r+1$ in $\mG$, and let
\[
        \mG,\mC,\mC_1,\ldots,\mC_i,
\]
be a divisor chain, each term being a proper divisor of the preceding one. Again there is no restriction in assuming that the chain begins with $\mG$. By 2$\gamma$ there is a composition series through $\mC$, and hence in $\mG/\mC$ there is a composition series of shorter length. By the induction hypothesis the divisor chain
\[
        \mG/\mC,\ \mC_1/\mC,\ldots,\mC_i/\mC,
\]
terminates after finitely many steps. By the one-to-one correspondence furnished by the first isomorphism theorem (§4, 2), the same holds for the original chain beginning with $\mG$, and hence also for the chain with $\mE$ adjoined. The equivalence of the two assumptions -- existence of a composition series and the double-chain condition -- is thereby proved.

\emph{Received August 13, 1925.}


\end{document}
